\hypertarget{native-291-remarks}{}\pdfbookmark[1]{Remarks on individual passages}{bemerkungen-remarks}
For page 45, line 1, there is a note by \textsc{Gauss} in his own copy saying that
$\omega'$ should stand there in place of $\omega$. However, $\omega$ is correct if
only the second term in equation XI, Art. 22, is calculated using hyperbolic logarithms
and the first ($e\operatorname{tang} F$) is calculated using Briggsian logarithms.
If \textls[150]{both} terms are calculated using hyperbolic logarithms, $\omega'$ must
be used and $\lambda$ must then be dropped from the denominator.

The following particular points should be noted about the numerical examples. Besides
Art. 150---155, the figures in Art. 10, 13, 14, 51, 63---65, 69, 77, 87\,I, 97\,I and
the calculation of the Sun's position in Art. 73 concern the planet \textls[150]{Juno}.
Besides Art. 156---157, Art. 56 and 58 and the calculation of the Sun's position in
Art. 70 concern the planet \textls[150]{Pallas}. Besides Art. 159, Art. 87\,II and 97\,II
concern the planet \textls[150]{Ceres}. The examples in Art. 38, 43, 97\,III also go
together for an ellipse of large eccentricity, and those in Art. 23, 24, 26, 46, 105
for a hyperbola.

On Art. 150---155. \textsc{Maskelyne} communicated the observations of Juno in a letter
to \textsc{Gauss} of 8 May 1805. \textsc{Gauss} took the positions of the Sun, the
nutation, and the obliquity of the ecliptic from \textsc{Zach}'s solar tables. A
detailed check of these values (page 195, lines 8---10) was dispensed with, as their
exact accuracy is immaterial to working through the examples and they could not have
been changed in view of the subsequent calculation. On page 195, line 8, however,
$53''{,}72$ was changed to $53''{,}71$ in the first position of the Sun to bring what
follows into agreement with it. The values of the constants that \textsc{Gauss} used
to reduce the observations are given below.
\[
\begin{aligned}
\text{Precession}&=50''{,}11\\
\text{Aberration}&=20''{,}00\\
\text{Paris-Greenwich longitude difference}&=9^{\mathrm m}20^{\mathrm s}{,}8.
\end{aligned}
\]

Some values of the elements (page 204) differed slightly when the calculation was
checked. The limit of precision with which they can be determined depends on how
precisely the quantity $x = \sin\tfrac12 g^2$ is known. Under the first hypothesis,
the \SourceLeaf{292}check gave
\[
\tfrac12 g = 1^0 34' 10''{,}046 \qquad x = 0{,}000\,7501\,446
\]
and from this the following values were obtained.
\[
\begin{aligned}
\text{Epoch} &= 42^0 12' 39''{,}89 \\
\text{Mean motion} &= 823''{,}4987 \\
\text{Perihelion} &= 52^0 41' 13''{,}24 \\
\varphi &= 14\ 24\ 28{,}58 \\
\text{Logarithm of the semimajor axis} &= 0{,}422\,8958 \\
\text{Ascending node} &= 171^0 5' 48''{,}84 \\
\text{Inclination} &= 13^0 2' 37''{,}50.
\end{aligned}
\]

The values given by \textsc{Gauss} seem to be based on the assumption
\[
\tfrac12 g = 1^0 34' 10''{,}056 \qquad x = 0{,}000\,7501\,474
\]
The differences thus fall below the attainable limit of precision. Under the second
hypothesis they were found to be even smaller. Under the third, \textsc{Gauss}'s
values were obtained exactly, except at the places noted in the list of errors.

On Art. 156---157. The declination of the third observation of Pallas is given in
\textsc{Zach}'s Monatlicher Correspondenz, at the place already cited, as
$-28^0 38' 8''{,}3$. The positions of the Sun on page 205, lines 5---7 from the bottom,
were left unchecked. \textsc{Gauss} presumably calculated the nutation using his own
tables, specifically the manuscript tables entered in a notebook (B\,b). Those tables
were preparatory work for the ones printed in the Monatlichen Correspondenz (see
Vol. VI, page 123) and differ from the printed ones only in the values of the constants.
The nutation was found to be correct except for the third value ($13''{,}68$ instead
of $12''{,}20$), and the aberration, calculated with the constant $20''{,}00$, except
for the third and fourth values ($21''{,}51$ and $1''{,}63$ instead of $23''{,}68$ and
$1''{,}31$). \textsc{Gauss} calculated the precession using the values $m = 45''{,}93$,
$n = 20''{,}02$. Here the third position differs ($2''{,}55$ instead of $2''{,}25$).
These differences could only partly be taken into account in the reprint (see the
list of errors).

Some of the elements on page 211 were also found to differ slightly, within the limits
of precision of the calculation as a whole. One example is the perihelion from the node,
$323^0 44' 56''{,}64$.

On Art. 159. The observation time for the first position of Ceres is given as
$13^{\mathrm h} 8^{\mathrm m} 47^{\mathrm s}$ in \textsc{Olbers}'s letter to \textsc{Gauss}
of 8 September 1805. The reduction is based on the differences in longitude Bremen-Paris
$= 25^{\mathrm m} 51^{\mathrm s}$, Göttingen-Paris $= 30^{\mathrm m} 21^{\mathrm s}$, and
Lilienthal-Paris $= 26^{\mathrm m} 14^{\mathrm s}$ given in \textsc{Zach}, Tabulae motuum
Solis, Supplementum, Gotha 1804. On page 213 the values for the obliquity of the ecliptic,
precession, and nutation in lines 8---11 and the positions of the Sun in lines 9 and 8
from the bottom were left unchecked. On page 214, more substantial inaccuracies were
found in the figures for hypothesis II, and only some of them could be corrected. The
values of $\omega+\sigma$, $\log Qc\sin\omega$, and $z$ were corrected. The value of $r'$
that follows then comes out as given by \textsc{Gauss}, so it was probably calculated
using the correct value of $z$ that has been put in. From here onward, however, the
calculation was correctly carried through using the incorrect value from the old edition
($z = 7^0 2' 32''{,}870$). On page 215 the calculation is correct as far as
$\log r' = 0{,}413\,2817$. From here, however, \textsc{Gauss} seems to have used the
incorrect values $\log\dfrac{n'r'}{n} = 0{,}662\,0412$ (instead of $0{,}662\,0380$) and
$\log\dfrac{n'r'}{n''} = 0{,}636\,4081$ (instead of $0{,}636\,4049$). The subsequent
calculation correctly follows from those values. When recalculated, the values for
the true anomalies, the angle $\varphi$, and the semimajor axis on page 215, lines
11, 10, 5, 4 from the bottom, also differed slightly, both with the
\SourceLeaf{293}
value $x = 0{,}065\,29078$ given in Art. 97 of the old edition and with the corrected
value $x = 0{,}065\,29076$. The values given by \textsc{Gauss} would correspond to
$g = 29^0 36' 31''{,}975$ ($x = 0{,}065\,290798$).

On Art. 171. For the first observation of \textls[150]{Vesta}, \textsc{Olbers} gives
the right ascension as $183^0 52' 37''$ in his letter to \textsc{Gauss} of 5 April 1807.
On page 231, the positions of the Sun in lines 12---15 and the derivation of the
figures in lines 9---6 from the bottom from the observations were left unchecked.
On page 235, the values of the elements appear to have been calculated with
$g = 23^0 34' 22''{,}34$ ($x = 0{,}041\,7240$). The figures on page 235, lines 7 to 2
from the bottom, were also left unchecked.

All the handwritten remarks that \textsc{Gauss} entered in his own copy of the
Theoria motus have been inserted as footnotes on pages 20, 30, 52, 53, 56, 65, 66,
80, 117---118, 120, 144, 145, 183, 226, 242, 249, 251, 270. In the note on page 118,
last line, $\dfrac{17420}{729}$ was corrected to $\dfrac{15420}{729} =
\dfrac{5120}{243}$, and on page 185, line 7 from the bottom,
$-3 \pm \BareRoot{(9-16\operatorname{cotang} A^2)}$ was corrected to
$-3\operatorname{cotang} A \pm \BareRoot{(9\operatorname{cotang} A^2-16)}$.
Expanding the expression in the footnote on p. 270 gives
$y = 1 + \dfrac{10}{9}h + \dfrac{110}{81}hh + \cdots$. A sign error was therefore
probably made (cf. p. 118), and a possible intended reading is
$y = \left(1+\dfrac{32}{9}h\right)^{\frac{5}{16}}$. For page 144, compare the remarks
on \textsc{Ivory}'s method for determining cometary orbits in \textsc{Gauss}'s letter
to \textsc{Olbers} of 31 December 1814.

The following further points may be noted.

On Art. 32\,IV. The maximum error in the true anomaly occurs at the value of $v$
determined by the equation
\begin{center}
$\cos v + e\cos 2v = 0$\quad or\quad \mbox{$\cos v = \dfrac1e\left\{ -\tfrac14 + \BareRoot{\left(\tfrac1{16}+\tfrac12 ee\right)}\right\}$}
.
\end{center}
For all the values of $e$ given on page 43, $v$ is approximately $60^0$.

On Art. 32\,VII. In the hyperbolic case, the maximum error in $F$ occurs when
$\cos F + \tfrac13\sin 2F - e\cos 2F = 0$. The values of $F$ corresponding to the
table on page 44 are given below.

\begin{center}
\begin{tabular}{c|c}
$e$ & $F$ \\
\hline
$1{,}3$ & $14^0 46'$ \\
$1{,}2$ & $11\ 16$ \\
$1{,}1$ & $6\ 41$ \\
$1{,}05$ & $3\ 44$ \\
$1{,}01$ & $0\ 50$ \\
$1{,}001$ & $0\ 5$ \\
\end{tabular}
\end{center}

On Art. 36. The maximum error in the elliptical case occurs when $1-2ee+e\cos E = 0$,
and in the hyperbolic case when $ee(1+6\sin F)+e(\cos F-3\operatorname{tang} F) = 2+3\sin F$.
Thus the values of $e$ corresponding to the tables on page 48 are given below.

\begin{center}
\begin{minipage}[t]{.42\linewidth}\centering
Ellipse\par\smallskip
\begin{tabular}{c|c}
$E$ & $e$ \\
\hline
$10^0$ & $0{,}99494$ \\
$20$ & $0{,}98003$ \\
$30$ & $0{,}95602$ \\
$40$ & $0{,}92409$ \\
$50$ & $0{,}88583$ \\
$60$ & $0{,}83534$ \\
\end{tabular}
\end{minipage}\hfill
\begin{minipage}[t]{.42\linewidth}\centering
Hyperbola\par\smallskip
\begin{tabular}{c|c}
$F$ & $e$ \\
\hline
$10^0$ & $1{,}00508$ \\
$20$ & $1{,}02094$ \\
$30$ & $1{,}04993$ \\
$40$ & $1{,}09754$ \\
$50$ & $1{,}17671$ \\
$60$ & $1{,}32007$ \\
$70$ & $1{,}63435$ \\
\end{tabular}
\end{minipage}
\end{center}
\SourceLeaf{294}
\textsc{Gauss}'s preparatory work for the Theoria motus is found chiefly in the notebooks
designated B\,b, B\,c, B\,d, and the tables in notebook B\,b, pp. 30, 31, 35. For the
composition of the manuscript and the preparations for printing, see the correspondence
with \textsc{Olbers} (W. \textsc{Olbers}, sein Leben und seine Werke, edited by C.
\textsc{Schilling}, Vol. II).

\begin{flushright}
\textsc{Brendel}.
\end{flushright}

\begin{center}
\rule{0.3\linewidth}{0.4pt}
\end{center}
